\chapter{Why compute with molecules?}
\label{ch:why-molecules}

\section*{The question behind the medium}
Imagine a delivery route through four locations. The route must start at $s$,
finish at $t$, and visit every location exactly once. A conventional program
can search for such a route. A molecular computer would instead represent
candidates in matter, transform the resulting population and interpret a
measurement. The interesting question is not whether the answer can be written
in DNA. It is whether the entire physical process implements the intended
computation, at an acceptable cost and with a defensible error model.

By the end of this chapter you should be able to distinguish a candidate from
a witness, prove what an ideal filter establishes, and explain why an empty
experimental readout is harder to interpret than an independently verified
positive result. No laboratory knowledge is assumed. Our first executable
example manipulates tokens, not molecules; making that boundary explicit is
part of learning the subject.

\section{A small problem with a revealing trap}
\begin{definition}[Directed Hamiltonian path with specified endpoints]
For a finite directed graph $G=(V,E)$ and designated vertices $s,t\in V$, a
witness is a sequence of $|V|$ distinct vertices beginning at $s$, ending at
$t$, and following an edge in $E$ at every consecutive step.
\end{definition}
The graph in our example has
\begin{align*}
 V &= \{s,u,v,t\},\\
 E &= \{(s,u),(s,v),(u,u),(u,v),(v,u),(u,t),(v,t)\}.
\end{align*}
An ordered pair $(u,v)$ permits travel from $u$ to $v$; it does not imply the
reverse direction. Here both directions happen to be present between $u$ and
$v$. The self-loop $(u,u)$ is intentional. It produces the walk $(s,u,u,t)$,
which has four entries and the correct endpoints but misses $v$. A test of
length is not a test of distinctness.

We use \emph{walk} for a sequence that follows the directed edges but may
repeat vertices. Length in this chapter counts vertex entries, not edges.
Thus $(s,u,t)$ has length three and uses two edges. Establishing this convention
before writing code prevents a surprisingly consequential off-by-one error.

There are two witnesses: $(s,u,v,t)$ and $(s,v,u,t)$. Knowing them in advance
lets us inspect the search process. It does not justify embedding these
answers in the implementation or declaring a filter correct merely because
it returns them once.

\section{Generate, filter, verify}
Our idealized algorithm first enumerates \emph{every} valid walk of lengths one
through $n=|V|$. It then retains walks with the designated endpoints, retains
those with exactly $n$ entries, and finally retains those containing every
vertex. Table~\ref{tab:pool} is generated by executing the reference code.
The initial population includes singleton walks and starts at every vertex;
it is not just the set of routes beginning at $s$.

\begin{table}[htbp]
\centering
% Generated by scripts/chapter01_artifacts.py; do not edit.
\begin{tabular}{lr}
\toprule Stage & Walks \\
\midrule
Generated & 45 \\
Endpoints & 5 \\
Length & 3 \\
Coverage & 2 \\
\bottomrule\end{tabular}

\caption{Exact token counts for the declared four-vertex graph. Counts refer
to distinct enumerated walks, not DNA molecules or experimental yields.}
\label{tab:pool}
\end{table}

The endpoint stage leaves $(s,u,t)$, $(s,v,t)$, $(s,u,u,t)$,
$(s,u,v,t)$ and $(s,v,u,t)$. Length removes the first two; coverage removes the
repeated-$u$ walk. These are different predicates with different jobs. Combining
them in an opaque function would hide the reason each is necessary.

\begin{proposition}[Soundness and conditional completeness]
Let the input population contain only valid walks of $G$. Every survivor of
the endpoint, length and coverage filters is a specified-endpoint Hamiltonian
path. If the population contains every valid walk of length $n$, every such
Hamiltonian path survives.
\end{proposition}
\begin{proof}
A survivor follows graph edges by the input precondition and has the required
endpoints by the first filter. It has $n$ positions by the second filter and
contains all $n$ vertices by the third. If any vertex appeared twice, at most
$n-1$ distinct vertices could occupy those $n$ positions, a contradiction.
Thus every vertex appears exactly once. Conversely, a Hamiltonian path of the
specified kind is present under the completeness assumption and satisfies
each predicate, so none of the filters removes it.
\end{proof}

The two halves have different assumptions. Soundness requires valid input
walks and correct predicates; completeness additionally requires the right
candidates to exist in the population. The proof does not say that a physical
process generates them, that a detector sees them or that searching is cheap.
It establishes a precisely bounded statement about finite sequences.

\section{Where molecules enter the story}
Adleman's 1994 experiment used DNA to solve a small directed Hamiltonian-path
instance with seven vertices. Its strategy generated molecular candidates,
selected for endpoint, length and vertex constraints, and read out a result.
The paper also discussed errors and the growth of material requirements
\citep{adleman1994}. Our four-vertex example is an original teaching model,
not a reconstruction of that experiment, its sequences or its measured yields.

What must change to turn our token computation into a molecular one? First,
there must be an encoding: physical states must represent vertices, adjacency
or candidate walks in a declared way. Next, actual transformations must
implement the predicates with known selectivity. Finally, a measurement must
be connected to an answer by a readout rule. A chemical product is not an answer
until we specify how to interpret it.

Figure~\ref{fig:pool-contract} separates the ideal search from these physical
obligations. The right-hand lane is not an extra filter that Python secretly
performs. It is the evidence that would be needed to justify a laboratory
realization. A sequence-based encoding alone does not establish such evidence.

\begin{figure}[p]
\centering
\includegraphics[width=\textwidth]{../build/figures/dna-g05.pdf}
\caption{DNA-G05. From exact search to experimental obligation. Solid arrows
carry candidates or evidence; dashed arrows identify a validation obligation.
The source TXT graph includes a reading guide and a failure condition. This
is a conceptual contract, not a measured experimental workflow.}
\label{fig:pool-contract}
\end{figure}

\section{An empty tube is not an impossibility proof}
Suppose a population contains only $(s,u,u,t)$. All filters behave exactly as
specified, yet no candidate survives. The graph still has two witnesses.
The population was incomplete. A negative result from this input therefore
does not establish that the problem has no solution.

This counterexample is intentionally nonchemical. It isolates a logical issue
before adding sampling, losses and measurement uncertainty. In a physical
realization, each of those processes could affect the relationship between
an empty readout and an empty set of mathematical solutions. The appropriate
response is not to assume failure, but to characterize that relationship with
controls and an explicit model.

For illustration, assume one particular witness is represented independently
in each of $M$ trials with probability $p$. The probability that it appears at
least once is $1-(1-p)^M$. This elementary expression follows by subtracting
the probability that all trials miss it. It is \emph{not} a prediction for a
DNA experiment: independence, $p$ and even what constitutes a trial must first
be justified. If $p=0$ because an encoding cannot produce the witness, increasing
$M$ never helps. A large population cannot repair a systematic omission.

A reported positive witness can be checked separately: confirm its endpoints,
count its entries, verify its set of vertices, and inspect every adjacent pair
against $E$. That check does not reproduce the whole search. It tests whether
the reported answer satisfies the problem, making it an especially useful
boundary between computation and interpretation.

\section{Parallel candidates still have a cost}
With $n$ available symbols there are $n^\ell$ words of length $\ell$ before edge
constraints. Our exhaustive population is bounded above by
\[
 \sum_{\ell=1}^{n} n^\ell.
\]
For the complete directed graph including self-loops, every such word is a
walk, so this bound is attained. Our sparse four-vertex graph has only 45
walks within the specified lengths. Neither number is a general statement
about the number of molecules needed by an experimental encoding.

Representing candidates simultaneously changes where work occurs; it does not
make representation free. Our program materializes tuples. An experiment
would need physical resources and readout capacity. Counting only three
filtering stages conceals the size of the population those stages act upon.
Counting only population size conceals the selectivity, time and measurement
cost of each stage. A useful resource account names both.

For the digital example, enumeration dominates the pedagogical design. This
is not presented as a competitive graph solver. Its virtue is that every
candidate and every rejection can be inspected. An efficient conventional
baseline would be necessary before making a performance claim.

\section{Reproduce and challenge the result}
After installing the repository's development dependencies, run the example:
\begin{Verbatim}[fontsize=\small]
from minidna.pool_search import chapter_example
result = chapter_example()
print(result.stages)
print(result.witnesses)
\end{Verbatim}
The artifact generator uses this same function to produce the table and a
machine-readable JSON record. The tests additionally compare all simple
directed graphs on three vertices, for each pair of distinct endpoints, with
an independent permutation-based oracle. That is 384 graph/endpoint cases,
not a proof for arbitrary $n$. The proposition provides the general argument;
the exhaustive small test checks implementation behavior over a declared domain.

The reference rejects malformed walks rather than allowing later predicates
to disguise an invalid edge. It also tests an incomplete population and the
single-vertex case. These checks are part of \gls{provenance}: the reader can
identify not only a result, but the assumptions and executable process that
produced it.\index{soundness}\index{completeness}\index{Hamiltonian path}

\section{Exercises and solution notes}
\begin{enumerate}
\item Remove $(u,u)$ from $E$. Predict all four stage counts before running the
program. Which logical distinction becomes less visible in this example?
\item Construct a walk that passes coverage but fails the length predicate.
Explain why coverage alone does not establish a Hamiltonian path.
\item Swap the order of length and coverage. Must the final result change?
Could intermediate population sizes change?
\item Supply an empty input population. State exactly what a zero-survivor
result proves and what it does not prove.
\item For $n=3$ with all nine directed edges including loops, compute the
number of generated walks of lengths one through three.
\item Design one independent witness check and one control for candidate
loss. Explain why the first cannot substitute for the second.
\end{enumerate}

\paragraph{Solution notes.}
(1) The counts become 26, 4, 2, 2; length no longer leaves a repeated-vertex
counterexample in this particular graph. (2) $(s,u,v,u,t)$ covers every vertex
but has five entries. (3) The final result is an intersection of predicates,
so order does not change it; intermediate sizes can differ. (4) No supplied
candidate survived; the graph may still have a witness. (5) $3+9+27=39$.
(6) An independent edge-and-vertex check establishes validity of an observed
witness; a known positive candidate processed through a proposed selection
workflow can probe loss under the control's conditions. Neither establishes
complete sampling of every possible witness.

\section*{What we have, and what comes next}
Next we introduce molecular structure and directional notation: prerequisites
for a DNA encoding, which our exact token model does not yet provide.
